% Curso: Programación en R — de cero a análisis de datos
% Compilar:  latexmk -pdf curso-r.tex      (o pdflatex dos veces)
% Los capítulos se generan desde los .md con:  python3 convertir.py
\documentclass[11pt, a4paper, oneside]{report}

% === Idioma y codificación ===
\usepackage[utf8]{inputenc}
\usepackage[T1]{fontenc}
\usepackage[spanish, es-tabla, es-nodecimaldot, es-noshorthands]{babel}
\usepackage{newunicodechar}
\newunicodechar{≥}{\ensuremath{\geq}}
\newunicodechar{≤}{\ensuremath{\leq}}
\newunicodechar{→}{\ensuremath{\rightarrow}}
\newunicodechar{²}{\textsuperscript{2}}

% === Tipografía ===
\usepackage{lmodern}
\usepackage{microtype}

% === Tablas (salida de pandoc) ===
\usepackage{longtable, booktabs, array, calc}
\renewcommand{\arraystretch}{1.25}

% === Layout y color ===
\usepackage[margin=2.5cm, headheight=14pt]{geometry}
\usepackage{xcolor}
\definecolor{rblue}{HTML}{1F4E78}
\definecolor{codebg}{HTML}{F5F7FA}
\definecolor{codegray}{HTML}{6A737D}
\definecolor{codestr}{HTML}{A31515}
\usepackage{fancyhdr}
\pagestyle{fancy}
\fancyhf{}
\fancyhead[L]{\small\nouppercase{\leftmark}}
\fancyhead[R]{\small Programación en R}
\fancyfoot[C]{\thepage}
\renewcommand{\chaptermark}[1]{\markboth{#1}{}}
\usepackage{titlesec}
\titleformat{\chapter}[display]
  {\normalfont\huge\bfseries\color{rblue}}{}{0pt}{\Huge}
\titlespacing*{\chapter}{0pt}{-20pt}{24pt}

% === Código R (listings; pandoc --listings) ===
\usepackage{listings}
\newcommand{\passthrough}[1]{#1}
\providecommand{\tightlist}{\setlength{\itemsep}{0pt}\setlength{\parskip}{0pt}}
\lstset{
  language=R,
  backgroundcolor=\color{codebg},
  basicstyle=\ttfamily\small,
  keywordstyle=\color{rblue}\bfseries,
  commentstyle=\color{codegray}\itshape,
  stringstyle=\color{codestr},
  breaklines=true,
  breakatwhitespace=false,
  postbreak=\mbox{\textcolor{codegray}{$\hookrightarrow$}\space},
  frame=single,
  rulecolor=\color{codegray!40},
  framesep=4pt,
  xleftmargin=6pt, xrightmargin=6pt,
  columns=fullflexible,
  keepspaces=true,
  showstringspaces=false,
  upquote=true,
  tabsize=2,
  % Caracteres UTF-8 dentro del código (pdflatex)
  literate=
    {á}{{\'a}}1 {é}{{\'e}}1 {í}{{\'i}}1 {ó}{{\'o}}1 {ú}{{\'u}}1
    {Á}{{\'A}}1 {É}{{\'E}}1 {Í}{{\'I}}1 {Ó}{{\'O}}1 {Ú}{{\'U}}1
    {ñ}{{\~n}}1 {Ñ}{{\~N}}1 {ü}{{\"u}}1 {¿}{{?`}}1 {¡}{{!`}}1
    {—}{{---}}1 {–}{{--}}1 {·}{{\textperiodcentered}}1
    {≥}{{$\geq$}}1 {≤}{{$\leq$}}1 {→}{{$\rightarrow$}}1 {²}{{\textsuperscript{2}}}1
}
% Quarto/Markdown no existen en listings: se muestran como texto plano.
\lstdefinelanguage{markdown}{}

% === Citas en bloque como notas destacadas ===
\usepackage{mdframed}
\renewenvironment{quote}
  {\begin{mdframed}[linecolor=rblue, linewidth=2pt, topline=false,
     bottomline=false, rightline=false, backgroundcolor=rblue!5,
     innertopmargin=6pt, innerbottommargin=6pt]}
  {\end{mdframed}}

% === Enlaces ===
\usepackage[colorlinks=true, linkcolor=rblue, urlcolor=rblue]{hyperref}
\usepackage{xurl}   % permite cortar URLs largas en cualquier carácter
\hypersetup{
  pdftitle={Programación en R: de cero a análisis de datos},
  pdfsubject={Curso modular de R},
  pdfkeywords={R, tidyverse, ggplot2, Quarto, análisis de datos}
}

\setcounter{tocdepth}{1}
\setlength{\parindent}{0pt}
\setlength{\parskip}{6pt plus 2pt minus 1pt}

\begin{document}

% === Portada ===
\hypersetup{pageanchor=false}
\begin{titlepage}
  \centering
  \vspace*{3cm}
  {\color{rblue}\rule{\linewidth}{1.5pt}\par}
  \vspace{0.8cm}
  {\Huge\bfseries\color{rblue} Programación en R\par}
  \vspace{0.4cm}
  {\LARGE de cero a análisis de datos\par}
  \vspace{0.8cm}
  {\color{rblue}\rule{\linewidth}{1.5pt}\par}
  \vspace{1.5cm}
  {\large Curso modular · 8 módulos · $\sim$34 horas\par}
  \vspace{0.3cm}
  {\large Nivel: principiante $\rightarrow$ intermedio\par}
  \vfill
  {\large Material didáctico con ejercicios, soluciones y rúbricas\par}
  \vspace{0.3cm}
  {\normalsize Código verificado con R 4.3.3 y tidyverse 2.0.0\par}
  \vspace{1cm}
  {\large \today\par}
\end{titlepage}
\hypersetup{pageanchor=true}

\tableofcontents

\chapter*{Presentación del curso}
\addcontentsline{toc}{chapter}{Presentación del curso}\markboth{Presentación del curso}{}

\textbf{Perfil:} capacitación en datos/analítica · \textbf{Nivel:}
principiante → intermedio · \textbf{Duración:} 8 módulos,
\textasciitilde34 h · \textbf{Modalidad:} autoformación o híbrida
(sesiones de 2 h + práctica) · \textbf{Idioma:} español

\section*{¿Para quién es?}
\addcontentsline{toc}{section}{¿Para quién es?}

Personas sin experiencia previa en programación (o que vienen de Excel)
que quieren usar R para limpiar, analizar, visualizar y reportar datos.
Los ejemplos usan un conjunto de \textbf{ventas simuladas de retail} (12
tiendas, 3 canales, 4 regiones, 5 productos, 52 semanas de 2025).

\section*{Objetivos del curso (diseño inverso)}
\addcontentsline{toc}{section}{Objetivos del curso (diseño inverso)}

Al terminar, la persona podrá:

\begin{enumerate}
\def\labelenumi{\arabic{enumi}.}
\tightlist
\item
  \textbf{Escribir} scripts de R legibles usando objetos, estructuras de
  datos, control de flujo y funciones propias.
\item
  \textbf{Importar y ordenar} datos de CSV/Excel siguiendo los
  principios de \emph{tidy data}.
\item
  \textbf{Transformar} y \textbf{resumir} datos con
  \passthrough{\lstinline!dplyr!} para calcular indicadores por grupo.
\item
  \textbf{Construir} visualizaciones claras con
  \passthrough{\lstinline!ggplot2!} adecuadas a cada pregunta.
\item
  \textbf{Aplicar e interpretar} estadística descriptiva, pruebas de
  hipótesis y regresión lineal.
\item
  \textbf{Producir} un reporte reproducible con Quarto que responda una
  pregunta de negocio.
\end{enumerate}

\textbf{Evidencia de logro:} ejercicios con solución en cada módulo, 3
entregas parciales (módulos 3, 5 y 7) y un \textbf{proyecto integrador}
evaluado con rúbrica (módulo 8).

\section*{Mapa del curso}
\addcontentsline{toc}{section}{Mapa del curso}

\begin{longtable}[]{@{}
  >{\raggedright\arraybackslash}p{(\columnwidth - 8\tabcolsep) * \real{0.2000}}
  >{\raggedright\arraybackslash}p{(\columnwidth - 8\tabcolsep) * \real{0.2000}}
  >{\raggedright\arraybackslash}p{(\columnwidth - 8\tabcolsep) * \real{0.2000}}
  >{\raggedright\arraybackslash}p{(\columnwidth - 8\tabcolsep) * \real{0.2000}}
  >{\raggedright\arraybackslash}p{(\columnwidth - 8\tabcolsep) * \real{0.2000}}@{}}
\toprule\noalign{}
\begin{minipage}[b]{\linewidth}\raggedright
\#
\end{minipage} & \begin{minipage}[b]{\linewidth}\raggedright
Módulo
\end{minipage} & \begin{minipage}[b]{\linewidth}\raggedright
Horas
\end{minipage} & \begin{minipage}[b]{\linewidth}\raggedright
Nivel de Bloom
\end{minipage} & \begin{minipage}[b]{\linewidth}\raggedright
Evaluación
\end{minipage} \\
\midrule\noalign{}
\endhead
\bottomrule\noalign{}
\endlastfoot
1 & Primeros pasos con R y RStudio & 4 & Recordar · Aplicar & Quiz \\
2 & Estructuras de datos & 4 & Comprender · Aplicar & Quiz + práctica \\
3 & Programación: control de flujo y funciones & 4 & Aplicar · Analizar
& \textbf{Entrega 1} (rúbrica) \\
4 & Importar y ordenar datos & 4 & Comprender · Aplicar & Quiz +
práctica \\
5 & Transformar datos con dplyr & 5 & Aplicar · Analizar &
\textbf{Entrega 2}: mini-reporte de KPIs \\
6 & Visualización con ggplot2 & 4 & Aplicar · Evaluar & Revisión entre
pares \\
7 & Estadística y modelos & 4 & Aplicar · Evaluar & \textbf{Entrega 3}:
informe breve \\
8 & Quarto y proyecto integrador & 5 & Crear & \textbf{Proyecto final}
(rúbrica global) \\
\end{longtable}

\section*{Ponderación sugerida}
\addcontentsline{toc}{section}{Ponderación sugerida}

\begin{longtable}[]{@{}ll@{}}
\toprule\noalign{}
Componente & Peso \\
\midrule\noalign{}
\endhead
\bottomrule\noalign{}
\endlastfoot
Quizzes y ejercicios (M1, M2, M4, M6) & 15 \% \\
Entrega 1 --- funciones (M3) & 15 \% \\
Entrega 2 --- KPIs con dplyr (M5) & 20 \% \\
Entrega 3 --- informe estadístico (M7) & 15 \% \\
Proyecto integrador (M8) & 35 \% \\
\end{longtable}

\section*{Preparación}
\addcontentsline{toc}{section}{Preparación}

\begin{enumerate}
\def\labelenumi{\arabic{enumi}.}
\item
  Instala \textbf{R} (\url{https://cran.r-project.org/}) y
  \textbf{RStudio Desktop}
  (\url{https://posit.co/download/rstudio-desktop/}).
\item
  Instala los paquetes del curso:

\begin{lstlisting}[language=R]
install.packages(c("tidyverse", "readxl", "broom", "scales", "renv"))
\end{lstlisting}
\item
  Abre esta carpeta como \textbf{Proyecto de RStudio} y genera los datos
  de práctica (ya incluidos en \passthrough{\lstinline!datos/!}):

\begin{lstlisting}[language=R]
source("datos/generar_ventas.R")
\end{lstlisting}
\end{enumerate}

Todo el código del curso se ejecutó y verificó con R 4.3.3 y tidyverse
2.0.0 usando este directorio como directorio de trabajo.

\section*{Estructura de archivos}
\addcontentsline{toc}{section}{Estructura de archivos}

\begin{lstlisting}
curso-r/
|-- README.md                          # este índice
|-- modulo-01-primeros-pasos.md
|-- modulo-02-estructuras-de-datos.md
|-- modulo-03-programacion.md
|-- modulo-04-importar-y-ordenar.md
|-- modulo-05-transformar-con-dplyr.md
|-- modulo-06-visualizacion-ggplot2.md
|-- modulo-07-estadistica-y-modelos.md
|-- modulo-08-reportes-quarto-y-proyecto.md
|-- datos/
|   |-- generar_ventas.R               # script reproducible (semilla fija)
|   |-- ventas.csv                     # 3 120 filas: fecha, tienda, producto, precio, unidades
|   `-- tiendas.csv                    # 12 tiendas: canal y región
`-- latex/
    |-- curso-r.tex                    # documento principal (portada, índice, preámbulo)
    |-- capitulos/                     # un .tex por módulo (generados)
    |-- convertir.py                   # regenera capitulos/ desde los .md (requiere pandoc)
    `-- curso-r.pdf                    # versión compilada
\end{lstlisting}

\textbf{Versión LaTeX/PDF:} desde \passthrough{\lstinline!latex/!},
ejecuta \passthrough{\lstinline!python3 convertir.py!} (solo si editaste
los \passthrough{\lstinline!.md!}) y luego
\passthrough{\lstinline!latexmk -pdf curso-r.tex!}. También se puede
subir la carpeta \passthrough{\lstinline!latex/!} a Overleaf y compilar
con pdfLaTeX.

\section*{Bibliografía general}
\addcontentsline{toc}{section}{Bibliografía general}

Leyenda: {[}ES{]} español · {[}EN{]} inglés · {[}OA{]} acceso abierto.

\subsection*{Libros}
\addcontentsline{toc}{subsection}{Libros}

\begin{itemize}
\tightlist
\item
  Grolemund, G. (2014). \emph{Hands-On Programming with R}. O'Reilly
  Media. ISBN 978-1-4493-5901-0.
  \url{https://rstudio-education.github.io/hopr/} {[}EN{]} {[}OA en
  línea{]}
\item
  Wickham, H. (2019). \emph{Advanced R} (2.ª ed.). CRC Press. ISBN
  978-0-8153-8457-1. \url{https://adv-r.hadley.nz/} {[}EN{]} {[}OA en
  línea{]}
\item
  Wickham, H., Çetinkaya-Rundel, M., \& Grolemund, G. (2023). \emph{R
  for Data Science: Import, tidy, transform, visualize, and model data}
  (2.ª ed.). O'Reilly Media. ISBN 978-1-4920-9740-2.
  \url{https://r4ds.hadley.nz/} {[}EN{]} {[}OA en línea{]}

  \begin{itemize}
  \tightlist
  \item
    Traducción al español (2.ª ed.):
    \url{https://davidrsch.github.io/r4ds-es/} {[}ES{]} {[}OA{]}
  \item
    Traducción comunitaria (1.ª ed.): \url{https://es.r4ds.hadley.nz/}
    {[}ES{]} {[}OA{]}
  \end{itemize}
\end{itemize}

\subsection*{Artículos}
\addcontentsline{toc}{subsection}{Artículos}

\begin{itemize}
\tightlist
\item
  Wickham, H. (2014). Tidy data. \emph{Journal of Statistical Software,
  59}(10), 1--23. \url{https://doi.org/10.18637/jss.v059.i10} {[}EN{]}
  {[}OA{]}
\item
  Wickham, H., Averick, M., Bryan, J., Chang, W., McGowan, L. D.,
  François, R., \ldots{} Yutani, H. (2019). Welcome to the tidyverse.
  \emph{Journal of Open Source Software, 4}(43), 1686.
  \url{https://doi.org/10.21105/joss.01686} {[}EN{]} {[}OA{]}
\end{itemize}

\subsection*{Documentación}
\addcontentsline{toc}{subsection}{Documentación}

\begin{itemize}
\tightlist
\item
  R Core Team. \emph{An Introduction to R}.
  \url{https://cran.r-project.org/doc/manuals/r-release/R-intro.html}
  {[}EN{]} {[}OA{]}
\item
  Quarto. \emph{Tutorial: Computations}.
  \url{https://quarto.org/docs/get-started/computations/rstudio}
  {[}EN{]}
\end{itemize}

\subsection*{Videos}
\addcontentsline{toc}{subsection}{Videos}

\begin{itemize}
\tightlist
\item
  freeCodeCamp.org. \emph{R Programming Tutorial -- Learn the Basics of
  Statistical Computing} {[}Video{]} (Barton Poulson, \textasciitilde2
  h). \url{https://youtu.be/_V8eKsto3Ug} {[}EN{]}
\item
  \emph{Curso completo de R para Data Science con Tidyverse} {[}Video{]}
  (Juan Gabriel Gomila Salas).
  \url{https://www.youtube.com/watch?v=9NJyhs5PlGQ} {[}ES{]}
  \emph{{[}canal y fecha POR VERIFICAR{]}}
\item
  StatQuest with Josh Starmer. \emph{Linear Regression in R,
  Step-by-Step} {[}Video{]}.
  \url{https://www.youtube.com/watch?v=u1cc1r_Y7M0} {[}EN{]}
\end{itemize}

\subsection*{Nota de verificación}
\addcontentsline{toc}{subsection}{Nota de verificación}

Las referencias se verificaron mediante búsqueda web (título, autores,
editorial, ISBN, DOI y existencia de la URL en índices de búsqueda) el
2026-09-25. No fue posible abrir directamente las páginas desde el
entorno de trabajo, por lo que las \textbf{fechas de publicación de los
videos} y el \textbf{nombre exacto del canal} del video en español
quedan marcados como \emph{POR VERIFICAR}. Revísalos antes de publicar
el curso.

\chapter{Módulo 1 --- Primeros pasos con R y RStudio}

\begin{itemize}
\tightlist
\item
  \textbf{Duración estimada:} 1.5 h teoría + 2.5 h práctica
\item
  \textbf{Objetivos de aprendizaje.} Al finalizar, la persona podrá:

  \begin{enumerate}
  \def\labelenumi{\arabic{enumi}.}
  \tightlist
  \item
    \textbf{Identificar} los paneles de RStudio (consola, editor,
    entorno, archivos/gráficos/ayuda).
  \item
    \textbf{Ejecutar} instrucciones en la consola y desde un script
    \passthrough{\lstinline!.R!}.
  \item
    \textbf{Crear} objetos con \passthrough{\lstinline!<-!} y
    \textbf{reconocer} sus tipos básicos
    (\passthrough{\lstinline!numeric!},
    \passthrough{\lstinline!integer!},
    \passthrough{\lstinline!character!},
    \passthrough{\lstinline!logical!}).
  \item
    \textbf{Usar} la ayuda (\passthrough{\lstinline!?!},
    \passthrough{\lstinline!help()!},
    \passthrough{\lstinline!example()!}) e \textbf{instalar} /
    \textbf{cargar} paquetes.
  \end{enumerate}
\end{itemize}

\section{Contenidos}

\begin{enumerate}
\def\labelenumi{\arabic{enumi}.}
\tightlist
\item
  Qué es R y qué es RStudio (el lenguaje vs.~el entorno de desarrollo).
\item
  Instalación: R desde CRAN y RStudio Desktop (Posit).
\item
  Proyectos de RStudio (\passthrough{\lstinline!.Rproj!}) y el
  directorio de trabajo.
\item
  La consola como calculadora; operadores aritméticos y de comparación.
\item
  Objetos y asignación; reglas de nombres
  (\passthrough{\lstinline!snake\_case!}).
\item
  Tipos de datos atómicos y \passthrough{\lstinline!class()!},
  \passthrough{\lstinline!typeof()!}.
\item
  Funciones: argumentos por posición y por nombre.
\item
  Paquetes: \passthrough{\lstinline!install.packages()!}
  vs.~\passthrough{\lstinline!library()!}.
\item
  Comentarios con \passthrough{\lstinline!\#!} y buenas prácticas de
  scripts.
\end{enumerate}

\section{Desarrollo}

\subsection{R como calculadora}

\begin{lstlisting}[language=R]
2 + 3 * 4        # precedencia: 14
(2 + 3) * 4      # 20
10 / 3           # 3.333333
10 %/% 3         # división entera: 3
10 %% 3          # residuo (módulo): 1
2^10             # potencia: 1024
sqrt(144)        # 12
\end{lstlisting}

\subsection{Objetos y asignación}

En R se asigna con \passthrough{\lstinline!<-!} (atajo en RStudio:
\passthrough{\lstinline!Alt!} + \passthrough{\lstinline!-!}). El nombre
va a la izquierda.

\begin{lstlisting}[language=R]
precio   <- 4999
unidades <- 12
ingreso  <- precio * unidades
ingreso            # 59988
\end{lstlisting}

\subsection{Tipos básicos}

\begin{lstlisting}[language=R]
class(3.14)        # "numeric"
class(5L)          # "integer" (la L fuerza entero)
class("Norte")     # "character"
class(TRUE)        # "logical"
is.numeric(precio) # TRUE
as.numeric("42")   # convierte texto a número: 42
\end{lstlisting}

\subsection{Funciones y argumentos}

\begin{lstlisting}[language=R]
round(3.14159, digits = 2)   # por nombre: 3.14
round(3.14159, 2)            # por posición: 3.14
seq(from = 1, to = 10, by = 3)
\end{lstlisting}

\subsection{Ayuda}

\begin{lstlisting}[language=R]
?mean              # abre la ayuda de mean()
help("seq")
example(mean)      # ejecuta los ejemplos de la ayuda
\end{lstlisting}

\subsection{Paquetes}

Un paquete se \textbf{instala una sola vez} por computadora y se
\textbf{carga en cada sesión}:

\begin{lstlisting}[language=R]
# install.packages("tidyverse")   # una vez (descomenta para instalar)
library(tidyverse)                # en cada sesión
\end{lstlisting}

\begin{quote}
\textbf{Buena práctica:} trabaja siempre dentro de un \textbf{Proyecto
de RStudio} (\passthrough{\lstinline!File → New Project!}). Así las
rutas son relativas (\passthrough{\lstinline!"datos/ventas.csv"!}) y
evitas \passthrough{\lstinline!setwd()!} con rutas absolutas que solo
funcionan en tu computadora.
\end{quote}

\section{Actividades y ejercicios}

\begin{enumerate}
\def\labelenumi{\arabic{enumi}.}
\tightlist
\item
  Crea un proyecto de RStudio llamado \passthrough{\lstinline!curso-r!}
  y un script \passthrough{\lstinline!modulo-01.R!}.
\item
  Calcula el ingreso de vender 35 audífonos de \$899 con 10 \% de
  descuento.
\item
  Guarda tu nombre en un objeto \passthrough{\lstinline!nombre!} y
  verifica su clase.
\item
  ¿Qué devuelve \passthrough{\lstinline!class(5 > 3)!}? ¿Y
  \passthrough{\lstinline!as.numeric(TRUE)!}?
\item
  Usa la ayuda para averiguar qué hace el argumento
  \passthrough{\lstinline!na.rm!} de \passthrough{\lstinline!mean()!}.
\end{enumerate}

\subsection{Soluciones}

\begin{lstlisting}[language=R]
# 2
ingreso_desc <- 35 * 899 * (1 - 0.10)
ingreso_desc                  # 28318.5

# 3
nombre <- "Ana"
class(nombre)                 # "character"

# 4
class(5 > 3)                  # "logical"
as.numeric(TRUE)              # 1

# 5: na.rm = TRUE elimina los NA antes de calcular
mean(c(10, NA, 20), na.rm = TRUE)   # 15
\end{lstlisting}

\section{Evaluación (formativa)}

Quiz de 5 preguntas:

\begin{enumerate}
\def\labelenumi{\arabic{enumi}.}
\tightlist
\item
  ¿Cuál es la diferencia entre R y RStudio?
\item
  ¿Qué operador asigna valores a un objeto en el estilo recomendado?
\item
  ¿Qué devuelve \passthrough{\lstinline!10 \%\% 4!}?
\item
  ¿Por qué \passthrough{\lstinline!install.packages()!} no se escribe en
  cada script?
\item
  Escribe una línea que calcule la raíz cuadrada de 2 redondeada a 3
  decimales.
\end{enumerate}

\emph{Respuestas:} 1) R es el lenguaje/intérprete; RStudio es el IDE. 2)
\passthrough{\lstinline!<-!}. 3) \passthrough{\lstinline!2!}. 4) Porque
solo se instala una vez; en el script basta
\passthrough{\lstinline!library()!}. 5)
\passthrough{\lstinline!round(sqrt(2), 3)!}.

\section{Referencias verificadas}

\begin{itemize}
\tightlist
\item
  \textbf{Web:} R Core Team. \emph{An Introduction to R}. CRAN.
  \url{https://cran.r-project.org/doc/manuals/r-release/R-intro.html}
  {[}EN{]} {[}OA{]}
\item
  \textbf{Libro:} Grolemund, G. (2014). \emph{Hands-On Programming with
  R}. O'Reilly Media. ISBN 978-1-4493-5901-0. Versión en línea:
  \url{https://rstudio-education.github.io/hopr/} {[}EN{]} {[}OA en
  línea{]} --- caps. 1--2.
\item
  \textbf{Libro:} Wickham, H., Çetinkaya-Rundel, M., \& Grolemund, G.
  (2023). \emph{R for Data Science} (2.ª ed.). O'Reilly Media. ISBN
  978-1-4920-9740-2. Traducción al español:
  \url{https://davidrsch.github.io/r4ds-es/} {[}ES{]} {[}OA{]} --- cap.
  ``Flujo de trabajo: conceptos básicos''.
\item
  \textbf{Video:} freeCodeCamp.org. \emph{R Programming Tutorial --
  Learn the Basics of Statistical Computing} {[}Video{]} (instructor:
  Barton Poulson). YouTube. \url{https://youtu.be/_V8eKsto3Ug} {[}EN{]}
  --- secciones ``Installing R'', ``RStudio'', ``Packages''.
  \emph{{[}fecha de publicación POR VERIFICAR{]}}
\end{itemize}

\chapter{Módulo 2 --- Estructuras de datos}

\begin{itemize}
\tightlist
\item
  \textbf{Duración estimada:} 1.5 h teoría + 2.5 h práctica
\item
  \textbf{Objetivos de aprendizaje.} Al finalizar, la persona podrá:

  \begin{enumerate}
  \def\labelenumi{\arabic{enumi}.}
  \tightlist
  \item
    \textbf{Crear} vectores, factores, listas, matrices y data frames.
  \item
    \textbf{Explicar} la coerción de tipos y el reciclaje de vectores.
  \item
    \textbf{Aplicar} indexación por posición, nombre y condición lógica
    (\passthrough{\lstinline![ ]!}, \passthrough{\lstinline![[ ]]!},
    \passthrough{\lstinline!$!}).
  \item
    \textbf{Clasificar} qué estructura conviene para un problema dado.
  \end{enumerate}
\end{itemize}

\section{Contenidos}

\begin{enumerate}
\def\labelenumi{\arabic{enumi}.}
\tightlist
\item
  Vectores atómicos: \passthrough{\lstinline!c()!},
  \passthrough{\lstinline!seq()!}, \passthrough{\lstinline!rep()!},
  \passthrough{\lstinline!length()!}.
\item
  Operaciones vectorizadas y reciclaje.
\item
  Coerción: jerarquía
  \passthrough{\lstinline!logical < integer < numeric < character!}.
\item
  Valores especiales: \passthrough{\lstinline!NA!},
  \passthrough{\lstinline!NULL!}, \passthrough{\lstinline!NaN!},
  \passthrough{\lstinline!Inf!}.
\item
  Factores: niveles y orden.
\item
  Matrices: \passthrough{\lstinline!matrix()!},
  \passthrough{\lstinline!dim()!}.
\item
  Listas: estructuras heterogéneas.
\item
  Data frames: la tabla de datos; \passthrough{\lstinline!str()!},
  \passthrough{\lstinline!head()!}, \passthrough{\lstinline!summary()!},
  \passthrough{\lstinline!nrow()!}.
\item
  Indexación: \passthrough{\lstinline![ ]!},
  \passthrough{\lstinline![[ ]]!}, \passthrough{\lstinline!$!} y filtros
  lógicos.
\end{enumerate}

\section{Desarrollo}

\subsection{Vectores y vectorización}

\begin{lstlisting}[language=R]
unidades <- c(12, 7, 20, 3, 15)
precios  <- c(4999, 7999, 899, 6499, 2499)
ingresos <- unidades * precios          # operación elemento a elemento
sum(ingresos)
mean(unidades)
length(unidades)
1:10 * 2                                # reciclaje de un escalar
\end{lstlisting}

\subsection{Coerción y NA}

\begin{lstlisting}[language=R]
c(1, "dos", TRUE)       # todo se vuelve character
c(1, TRUE, FALSE)       # 1 1 0 (logical -> numeric)
x <- c(4, NA, 10)
sum(x)                  # NA
sum(x, na.rm = TRUE)    # 14
is.na(x)                # FALSE TRUE FALSE
\end{lstlisting}

\subsection{Factores}

\begin{lstlisting}[language=R]
canal <- factor(c("En línea", "Departamental", "En línea", "Autoservicio"))
levels(canal)                   # orden alfabético por defecto
table(canal)
talla <- factor(c("M", "S", "L", "M"), levels = c("S", "M", "L"), ordered = TRUE)
talla < "L"
\end{lstlisting}

\subsection{Matrices y listas}

\begin{lstlisting}[language=R]
m <- matrix(1:6, nrow = 2)
dim(m)
m[2, 3]                         # fila 2, columna 3

tienda <- list(id = "T01", region = "Norte", ventas = c(120, 98, 143))
tienda$region
tienda[["ventas"]][2]           # 98
str(tienda)
\end{lstlisting}

\subsection{Data frames}

\begin{lstlisting}[language=R]
df <- data.frame(
  producto = c("Smartphone A", "Tablet", "Audífonos"),
  precio   = c(4999, 6499, 899),
  unidades = c(12, 4, 20)
)
str(df)
df$ingreso <- df$precio * df$unidades   # nueva columna
df[df$unidades > 10, ]                  # filas que cumplen condición
df[, c("producto", "ingreso")]          # columnas por nombre
nrow(df); ncol(df)
\end{lstlisting}

\subsection{Leer el conjunto de datos del curso (R base)}

\begin{lstlisting}[language=R]
ventas <- read.csv("datos/ventas.csv")
head(ventas)
summary(ventas$unidades)
\end{lstlisting}

\section{Actividades y ejercicios}

\begin{enumerate}
\def\labelenumi{\arabic{enumi}.}
\tightlist
\item
  Crea un vector con las temperaturas
  \passthrough{\lstinline!c(22, 25, NA, 30, 28)!} y calcula su promedio
  ignorando el \passthrough{\lstinline!NA!}.
\item
  Con \passthrough{\lstinline!ventas!}, ¿cuántas filas y columnas tiene?
  ¿Cuántos \passthrough{\lstinline!NA!} hay en
  \passthrough{\lstinline!unidades!}?
\item
  Convierte \passthrough{\lstinline!ventas$producto!} en factor y
  muestra cuántos registros hay por producto.
\item
  Extrae las filas de \passthrough{\lstinline!ventas!} del producto
  \passthrough{\lstinline!"Tablet"!} con más de 8 unidades.
\item
  Crea una lista \passthrough{\lstinline!resumen!} con el total de
  unidades y el producto más frecuente.
\end{enumerate}

\subsection{Soluciones}

\begin{lstlisting}[language=R]
# 1
temp <- c(22, 25, NA, 30, 28)
mean(temp, na.rm = TRUE)                       # 26.25

# 2
dim(ventas)
sum(is.na(ventas$unidades))                    # 40

# 3
ventas$producto <- factor(ventas$producto)
table(ventas$producto)

# 4
tablets <- ventas[ventas$producto == "Tablet" &
                  !is.na(ventas$unidades) &
                  ventas$unidades > 8, ]
head(tablets)

# 5
resumen <- list(
  total_unidades = sum(ventas$unidades, na.rm = TRUE),
  mas_frecuente  = names(which.max(table(ventas$producto)))
)
str(resumen)
\end{lstlisting}

\section{Evaluación (formativa)}

\begin{itemize}
\tightlist
\item
  \textbf{Quiz:} ¿Qué devuelve \passthrough{\lstinline!c(1, "a")!}?
  ¿Cuál es la diferencia entre \passthrough{\lstinline!lista["x"]!} y
  \passthrough{\lstinline!lista[["x"]]!}? ¿Para qué sirve
  \passthrough{\lstinline!levels=!} en
  \passthrough{\lstinline!factor()!}?
\item
  \textbf{Práctica calificada:} script que lea
  \passthrough{\lstinline!ventas.csv!}, cree la columna
  \passthrough{\lstinline!ingreso!} y muestre las 5 filas con mayor
  ingreso (pista: \passthrough{\lstinline!order()!}, con
  \passthrough{\lstinline!decreasing = TRUE!}).
\end{itemize}

\section{Referencias verificadas}

\begin{itemize}
\tightlist
\item
  \textbf{Libro:} Grolemund, G. (2014). \emph{Hands-On Programming with
  R}. O'Reilly Media. ISBN 978-1-4493-5901-0.
  \url{https://rstudio-education.github.io/hopr/} {[}EN{]} {[}OA en
  línea{]} --- cap. ``R Objects'' y ``R Notation''.
\item
  \textbf{Libro:} Wickham, H. (2019). \emph{Advanced R} (2.ª ed.). CRC
  Press. ISBN 978-0-8153-8457-1. \url{https://adv-r.hadley.nz/} {[}EN{]}
  {[}OA en línea, CC BY-NC-SA 4.0{]} --- cap. 3 ``Vectors'' y cap. 4
  ``Subsetting'' (lectura de profundización).
\item
  \textbf{Web:} R Core Team. \emph{An Introduction to R} --- secciones
  ``Simple manipulations; numbers and vectors'' y ``Lists and data
  frames''.
  \url{https://cran.r-project.org/doc/manuals/r-release/R-intro.html}
  {[}EN{]} {[}OA{]}
\item
  \textbf{Video:} freeCodeCamp.org. \emph{R Programming Tutorial --
  Learn the Basics of Statistical Computing} {[}Video{]}. YouTube.
  \url{https://youtu.be/_V8eKsto3Ug} {[}EN{]} --- secciones ``Data
  Formats'', ``Factors'', ``Entering Data'', ``Importing Data''.
\end{itemize}

\chapter{Módulo 3 --- Programación: control de flujo y funciones}

\begin{itemize}
\tightlist
\item
  \textbf{Duración estimada:} 1.5 h teoría + 2.5 h práctica
\item
  \textbf{Objetivos de aprendizaje.} Al finalizar, la persona podrá:

  \begin{enumerate}
  \def\labelenumi{\arabic{enumi}.}
  \tightlist
  \item
    \textbf{Implementar} decisiones con
    \passthrough{\lstinline!if!}/\passthrough{\lstinline!else!},
    \passthrough{\lstinline!ifelse()!} y
    \passthrough{\lstinline!dplyr::case\_when()!}.
  \item
    \textbf{Ejecutar} iteraciones con \passthrough{\lstinline!for!} y
    \passthrough{\lstinline!while!}.
  \item
    \textbf{Escribir} funciones propias con argumentos, valores por
    defecto y validaciones.
  \item
    \textbf{Comparar} un bucle con la alternativa vectorizada o
    funcional (\passthrough{\lstinline!sapply()!},
    \passthrough{\lstinline!purrr::map()!}).
  \end{enumerate}
\end{itemize}

\section{Contenidos}

\begin{enumerate}
\def\labelenumi{\arabic{enumi}.}
\tightlist
\item
  Operadores lógicos: \passthrough{\lstinline!\&!},
  \passthrough{\lstinline!|!}, \passthrough{\lstinline"!"},
  \passthrough{\lstinline!\&\&!}, \passthrough{\lstinline!||!},
  \passthrough{\lstinline!\%in\%!}.
\item
  \passthrough{\lstinline!if!}, \passthrough{\lstinline!else if!},
  \passthrough{\lstinline!else!}; \passthrough{\lstinline!ifelse()!}
  vectorizado.
\item
  Bucles \passthrough{\lstinline!for!} y
  \passthrough{\lstinline!while!}; \passthrough{\lstinline!break!} y
  \passthrough{\lstinline!next!}.
\item
  Anatomía de una función: nombre, argumentos, cuerpo, valor de retorno.
\item
  Valores por defecto y validación con \passthrough{\lstinline!stop()!}.
\item
  Alcance (scope) de variables.
\item
  Vectorización y familia \passthrough{\lstinline!apply!}
  (\passthrough{\lstinline!sapply!}, \passthrough{\lstinline!lapply!})
  vs.~\passthrough{\lstinline!purrr::map()!}.
\end{enumerate}

\section{Desarrollo}

\subsection{Condicionales}

\begin{lstlisting}[language=R]
unidades <- 25
if (unidades >= 20) {
  "Alta rotación"
} else if (unidades >= 10) {
  "Rotación media"
} else {
  "Baja rotación"
}

u <- c(3, 12, 25)
ifelse(u >= 10, "Suficiente", "Reabastecer")   # vectorizado
"Tablet" %in% c("Tablet", "Reloj")             # TRUE
\end{lstlisting}

\subsection{Bucles}

\begin{lstlisting}[language=R]
productos <- c("Smartphone A", "Tablet", "Reloj")
for (p in productos) {
  print(paste("Revisando inventario de", p))
}

stock <- 50
semanas <- 0
while (stock > 0) {
  stock <- stock - 12
  semanas <- semanas + 1
}
semanas        # semanas hasta agotar: 5
\end{lstlisting}

\subsection{Funciones}

\begin{lstlisting}[language=R]
ingreso <- function(precio, unidades, descuento = 0) {
  if (descuento < 0 || descuento >= 1) {
    stop("`descuento` debe estar en [0, 1).")
  }
  precio * unidades * (1 - descuento)
}

ingreso(4999, 10)
ingreso(4999, 10, descuento = 0.15)
ingreso(c(899, 2499), c(20, 8))     # funciona con vectores
\end{lstlisting}

\subsection{Vectorizar en lugar de iterar}

\begin{lstlisting}[language=R]
x <- 1:5
# Con bucle
cuadrados <- numeric(length(x))
for (i in seq_along(x)) cuadrados[i] <- x[i]^2
cuadrados
# Vectorizado (preferido)
x^2
# Funcional
sapply(x, function(v) v^2)
library(purrr)
map_dbl(x, \(v) v^2)               # \(v) es la sintaxis corta de function(v), R >= 4.1
\end{lstlisting}

\section{Actividades y ejercicios}

\begin{enumerate}
\def\labelenumi{\arabic{enumi}.}
\tightlist
\item
  Escribe \passthrough{\lstinline!clasificar\_rotacion(u)!} que devuelva
  \passthrough{\lstinline!"Alta"!}, \passthrough{\lstinline!"Media"!} o
  \passthrough{\lstinline!"Baja"!} para un \textbf{vector} de unidades
  (umbrales 20 y 10).
\item
  Escribe \passthrough{\lstinline!margen(precio, costo)!} que devuelva
  el margen porcentual y lance error si
  \passthrough{\lstinline!costo > precio!}.
\item
  Con un bucle \passthrough{\lstinline!for!}, imprime el total de
  unidades por producto en \passthrough{\lstinline!ventas!}. Luego obtén
  lo mismo con \passthrough{\lstinline!tapply()!}.
\item
  Simula lanzar un dado 1 000 veces con
  \passthrough{\lstinline!sample()!} y calcula la proporción de seises.
\end{enumerate}

\subsection{Soluciones}

\begin{lstlisting}[language=R]
# 1 (dplyr::case_when es vectorizado y legible)
clasificar_rotacion <- function(u) {
  dplyr::case_when(
    u >= 20 ~ "Alta",
    u >= 10 ~ "Media",
    TRUE    ~ "Baja"
  )
}
clasificar_rotacion(c(3, 12, 25))

# 2
margen <- function(precio, costo) {
  if (any(costo > precio)) stop("El costo no puede ser mayor que el precio.")
  (precio - costo) / precio * 100
}
margen(100, 60)                   # 40

# 3
ventas <- read.csv("datos/ventas.csv")
for (p in unique(ventas$producto)) {
  total <- sum(ventas$unidades[ventas$producto == p], na.rm = TRUE)
  cat(p, ":", total, "\n")
}
tapply(ventas$unidades, ventas$producto, sum, na.rm = TRUE)

# 4
set.seed(1)
dado <- sample(1:6, 1000, replace = TRUE)
mean(dado == 6)
\end{lstlisting}

\section{Evaluación (sumativa parcial)}

Entrega un script \passthrough{\lstinline!funciones\_ventas.R!} con 3
funciones documentadas (comentario de propósito, argumentos y retorno) y
al menos una validación con \passthrough{\lstinline!stop()!}. Rúbrica:

\begin{longtable}[]{@{}
  >{\raggedright\arraybackslash}p{(\columnwidth - 6\tabcolsep) * \real{0.2500}}
  >{\raggedright\arraybackslash}p{(\columnwidth - 6\tabcolsep) * \real{0.2500}}
  >{\raggedright\arraybackslash}p{(\columnwidth - 6\tabcolsep) * \real{0.2500}}
  >{\raggedright\arraybackslash}p{(\columnwidth - 6\tabcolsep) * \real{0.2500}}@{}}
\toprule\noalign{}
\begin{minipage}[b]{\linewidth}\raggedright
Criterio
\end{minipage} & \begin{minipage}[b]{\linewidth}\raggedright
3 --- Logrado
\end{minipage} & \begin{minipage}[b]{\linewidth}\raggedright
2 --- En proceso
\end{minipage} & \begin{minipage}[b]{\linewidth}\raggedright
1 --- Inicial
\end{minipage} \\
\midrule\noalign{}
\endhead
\bottomrule\noalign{}
\endlastfoot
Correctitud & Las 3 funciones dan el resultado esperado & 2 funciones
correctas & ≤ 1 correcta \\
Vectorización & Aceptan vectores sin bucles innecesarios & Parcial &
Solo escalares \\
Validación & Mensajes de error claros & Validación incompleta & Sin
validación \\
Estilo & Nombres claros, comentarios útiles & Aceptable & Difícil de
leer \\
\end{longtable}

\section{Referencias verificadas}

\begin{itemize}
\tightlist
\item
  \textbf{Libro:} Wickham, H., Çetinkaya-Rundel, M., \& Grolemund, G.
  (2023). \emph{R for Data Science} (2.ª ed.). O'Reilly Media. ISBN
  978-1-4920-9740-2. Español: \url{https://davidrsch.github.io/r4ds-es/}
  {[}ES{]} {[}OA{]} --- caps. ``Funciones'' e ``Iteración''.
\item
  \textbf{Libro:} Grolemund, G. (2014). \emph{Hands-On Programming with
  R}. O'Reilly Media. ISBN 978-1-4493-5901-0.
  \url{https://rstudio-education.github.io/hopr/} {[}EN{]} {[}OA en
  línea{]} --- proyecto ``Slot machine'' (if/else, bucles, funciones).
\item
  \textbf{Libro:} Wickham, H. (2019). \emph{Advanced R} (2.ª ed.). CRC
  Press. ISBN 978-0-8153-8457-1. \url{https://adv-r.hadley.nz/} {[}EN{]}
  {[}OA en línea{]} --- cap. 5 ``Control flow'' y cap. 6 ``Functions''.
\end{itemize}

\chapter{Módulo 4 --- Importar y ordenar datos (tidy data)}

\begin{itemize}
\tightlist
\item
  \textbf{Duración estimada:} 1.5 h teoría + 2.5 h práctica
\item
  \textbf{Objetivos de aprendizaje.} Al finalizar, la persona podrá:

  \begin{enumerate}
  \def\labelenumi{\arabic{enumi}.}
  \tightlist
  \item
    \textbf{Explicar} los tres principios de los datos ordenados
    (\emph{tidy data}).
  \item
    \textbf{Importar} archivos CSV y Excel con
    \passthrough{\lstinline!readr!} y \passthrough{\lstinline!readxl!},
    controlando tipos de columna.
  \item
    \textbf{Transformar} tablas entre formato ancho y largo con
    \passthrough{\lstinline!pivot\_longer()!} /
    \passthrough{\lstinline!pivot\_wider()!}.
  \item
    \textbf{Diferenciar} los problemas de calidad más comunes (tipos
    erróneos, \passthrough{\lstinline!NA!}, duplicados).
  \end{enumerate}
\end{itemize}

\section{Contenidos}

\begin{enumerate}
\def\labelenumi{\arabic{enumi}.}
\tightlist
\item
  El \emph{tidyverse} y los \emph{tibbles}.
\item
  Principios de datos ordenados: cada variable una columna, cada
  observación una fila, cada valor una celda.
\item
  \passthrough{\lstinline!readr::read\_csv()!}, especificación de tipos
  con \passthrough{\lstinline!col\_types!},
  \passthrough{\lstinline!problems()!}.
\item
  \passthrough{\lstinline!readxl::read\_excel()!}: hojas y rangos.
\item
  \passthrough{\lstinline!tidyr::pivot\_longer()!} y
  \passthrough{\lstinline!pivot\_wider()!}.
\item
  \passthrough{\lstinline!separate\_wider\_delim()!} y
  \passthrough{\lstinline!unite()!}.
\item
  Manejo de \passthrough{\lstinline!NA!}:
  \passthrough{\lstinline!drop\_na()!},
  \passthrough{\lstinline!replace\_na()!}; duplicados con
  \passthrough{\lstinline!distinct()!}.
\item
  Exportar: \passthrough{\lstinline!write\_csv()!}.
\end{enumerate}

\section{Desarrollo}

\subsection{Importar con readr}

\begin{lstlisting}[language=R]
library(tidyverse)

ventas <- read_csv("datos/ventas.csv", col_types = cols(
  fecha     = col_date(),
  tienda_id = col_character(),
  producto  = col_character(),
  precio    = col_double(),
  unidades  = col_integer()
))
ventas
glimpse(ventas)
tiendas <- read_csv("datos/tiendas.csv", show_col_types = FALSE)
\end{lstlisting}

\subsection{Excel}

\begin{lstlisting}[language=R]
library(readxl)
# excel_sheets("datos/archivo.xlsx")                 # lista las hojas
# read_excel("datos/archivo.xlsx", sheet = "Ventas", range = "A1:E500")
\end{lstlisting}

\subsection{De ancho a largo y de regreso}

Muchas hojas de cálculo guardan un mes por columna: eso \textbf{no} es
\emph{tidy}, porque ``mes'' es una variable.

\begin{lstlisting}[language=R]
reporte_ancho <- tribble(
  ~tienda, ~ene, ~feb, ~mar,
  "T01",    120,   98,  143,
  "T02",     87,  110,   95
)

reporte_largo <- reporte_ancho |>
  pivot_longer(cols = ene:mar, names_to = "mes", values_to = "unidades")
reporte_largo

reporte_largo |>
  pivot_wider(names_from = mes, values_from = unidades)
\end{lstlisting}

\begin{quote}
\passthrough{\lstinline!|>!} es el \emph{pipe} nativo de R (≥ 4.1):
\passthrough{\lstinline!x |> f(y)!} equivale a
\passthrough{\lstinline!f(x, y)!}. Léelo como ``y luego''.
\end{quote}

\subsection{Separar y unir columnas}

\begin{lstlisting}[language=R]
codigos <- tibble(sku = c("CEL-001-NEGRO", "TAB-014-GRIS"))
codigos |>
  separate_wider_delim(sku, delim = "-", names = c("categoria", "numero", "color"))
\end{lstlisting}

\subsection{Valores faltantes y duplicados}

\begin{lstlisting}[language=R]
ventas |> summarise(faltantes = sum(is.na(unidades)))
ventas_sin_na <- ventas |> drop_na(unidades)
ventas |> mutate(unidades = replace_na(unidades, 0L)) |> head()
ventas |> distinct(tienda_id, fecha, producto) |> nrow()   # verifica llave única
\end{lstlisting}

\section{Actividades y ejercicios}

\begin{enumerate}
\def\labelenumi{\arabic{enumi}.}
\tightlist
\item
  Importa \passthrough{\lstinline!ventas.csv!} con
  \passthrough{\lstinline!read\_csv()!} sin
  \passthrough{\lstinline!col\_types!} y compara el tipo que asigna a
  \passthrough{\lstinline!unidades!} y \passthrough{\lstinline!fecha!}
  con la versión especificada.
\item
  Convierte \passthrough{\lstinline!ventas!} a formato ancho: una fila
  por \passthrough{\lstinline!fecha!} y
  \passthrough{\lstinline!tienda\_id!}, una columna por producto con las
  unidades.
\item
  Regresa el resultado anterior a formato largo y comprueba que
  recuperas el mismo número de filas.
\item
  Cuenta los \passthrough{\lstinline!NA!} de
  \passthrough{\lstinline!unidades!} por producto.
\end{enumerate}

\subsection{Soluciones}

\begin{lstlisting}[language=R]
# 1
v2 <- read_csv("datos/ventas.csv", show_col_types = FALSE)
spec(v2)          # readr infiere fecha como date y unidades como double

# 2
ancho <- ventas |>
  select(fecha, tienda_id, producto, unidades) |>
  pivot_wider(names_from = producto, values_from = unidades)
ancho

# 3
largo <- ancho |>
  pivot_longer(-c(fecha, tienda_id), names_to = "producto", values_to = "unidades")
nrow(largo) == nrow(ventas)    # TRUE

# 4
ventas |>
  group_by(producto) |>
  summarise(n_na = sum(is.na(unidades)))
\end{lstlisting}

\section{Evaluación (formativa)}

\begin{itemize}
\tightlist
\item
  \textbf{Quiz:} Da un ejemplo de tabla que viole cada principio de
  \emph{tidy data}. ¿Qué función usarías para arreglarla?
\item
  \textbf{Práctica:} Recibe un Excel ``reporte mensual'' con meses en
  columnas y entrégalo en formato largo como CSV con
  \passthrough{\lstinline!write\_csv()!}.
\end{itemize}

\section{Referencias verificadas}

\begin{itemize}
\tightlist
\item
  \textbf{Artículo:} Wickham, H. (2014). Tidy data. \emph{Journal of
  Statistical Software, 59}(10), 1--23.
  \url{https://doi.org/10.18637/jss.v059.i10} {[}EN{]} {[}OA{]}
\item
  \textbf{Artículo:} Wickham, H., Averick, M., Bryan, J., Chang, W.,
  McGowan, L. D., François, R., \ldots{} Yutani, H. (2019). Welcome to
  the tidyverse. \emph{Journal of Open Source Software, 4}(43), 1686.
  \url{https://doi.org/10.21105/joss.01686} {[}EN{]} {[}OA, CC BY 4.0{]}
\item
  \textbf{Libro:} Wickham, H., Çetinkaya-Rundel, M., \& Grolemund, G.
  (2023). \emph{R for Data Science} (2.ª ed.). O'Reilly Media. ISBN
  978-1-4920-9740-2. Español: \url{https://davidrsch.github.io/r4ds-es/}
  {[}ES{]} {[}OA{]} --- caps. ``Ordenar datos'', ``Importación de
  datos'', ``Hojas de cálculo''.
\item
  \textbf{Video:} Gomila Salas, J. G. (Frogames). \emph{Curso completo
  de R para Data Science con Tidyverse -- R for Data Science en Español}
  {[}Video{]}. YouTube.
  \url{https://www.youtube.com/watch?v=9NJyhs5PlGQ} {[}ES{]}
  \emph{{[}nombre exacto del canal y fecha POR VERIFICAR{]}}
\end{itemize}

\chapter{Módulo 5 --- Transformar datos con dplyr}

\begin{itemize}
\tightlist
\item
  \textbf{Duración estimada:} 1.5 h teoría + 3.5 h práctica
\item
  \textbf{Objetivos de aprendizaje.} Al finalizar, la persona podrá:

  \begin{enumerate}
  \def\labelenumi{\arabic{enumi}.}
  \tightlist
  \item
    \textbf{Aplicar} los verbos de \passthrough{\lstinline!dplyr!}
    (\passthrough{\lstinline!filter!}, \passthrough{\lstinline!select!},
    \passthrough{\lstinline!mutate!}, \passthrough{\lstinline!arrange!},
    \passthrough{\lstinline!summarise!}) encadenados con
    \passthrough{\lstinline!|>!}.
  \item
    \textbf{Calcular} indicadores agregados por grupo con
    \passthrough{\lstinline!group\_by()!} /
    \passthrough{\lstinline!.by!}.
  \item
    \textbf{Combinar} tablas con uniones
    (\passthrough{\lstinline!left\_join!},
    \passthrough{\lstinline!inner\_join!},
    \passthrough{\lstinline!anti\_join!}).
  \item
    \textbf{Usar} \passthrough{\lstinline!lubridate!} y
    \passthrough{\lstinline!stringr!} para trabajar con fechas y texto.
  \item
    \textbf{Analizar} el desempeño por canal, región y producto a partir
    de datos crudos.
  \end{enumerate}
\end{itemize}

\section{Contenidos}

\begin{enumerate}
\def\labelenumi{\arabic{enumi}.}
\tightlist
\item
  Los cinco verbos principales y el \emph{pipe}.
\item
  \passthrough{\lstinline!mutate()!} con
  \passthrough{\lstinline!if\_else()!} y
  \passthrough{\lstinline!case\_when()!}.
\item
  \passthrough{\lstinline!group\_by()!} +
  \passthrough{\lstinline!summarise()!};
  \passthrough{\lstinline!count()!}; agrupación temporal con
  \passthrough{\lstinline!.by!}.
\item
  Funciones de ventana: \passthrough{\lstinline!lag()!},
  \passthrough{\lstinline!cumsum()!}, \passthrough{\lstinline!rank()!},
  \passthrough{\lstinline!slice\_max()!}.
\item
  Uniones de tablas y llaves.
\item
  Fechas con \passthrough{\lstinline!lubridate!}:
  \passthrough{\lstinline!year()!}, \passthrough{\lstinline!month()!},
  \passthrough{\lstinline!floor\_date()!}.
\item
  Texto con \passthrough{\lstinline!stringr!}:
  \passthrough{\lstinline!str\_detect()!},
  \passthrough{\lstinline!str\_to\_upper()!},
  \passthrough{\lstinline!str\_replace()!}.
\end{enumerate}

\section{Desarrollo}

\begin{lstlisting}[language=R]
library(tidyverse)
ventas  <- read_csv("datos/ventas.csv", show_col_types = FALSE)
tiendas <- read_csv("datos/tiendas.csv", show_col_types = FALSE)
\end{lstlisting}

\subsection{Los verbos básicos}

\begin{lstlisting}[language=R]
ventas |>
  filter(producto == "Tablet", unidades >= 8) |>
  select(fecha, tienda_id, unidades) |>
  arrange(desc(unidades))

ventas <- ventas |>
  mutate(ingreso = precio * unidades)
\end{lstlisting}

\subsection{Resúmenes por grupo}

\begin{lstlisting}[language=R]
ventas |>
  group_by(producto) |>
  summarise(
    unidades = sum(unidades, na.rm = TRUE),
    ingreso  = sum(ingreso,  na.rm = TRUE),
    ticket_promedio = ingreso / unidades
  ) |>
  arrange(desc(ingreso))

ventas |> count(producto)                         # conteo rápido
\end{lstlisting}

\subsection{Uniones}

\begin{lstlisting}[language=R]
ventas_det <- ventas |>
  left_join(tiendas, by = "tienda_id")

ventas_det |>
  summarise(ingreso = sum(ingreso, na.rm = TRUE), .by = c(canal, region)) |>
  arrange(canal, desc(ingreso))

# ¿Hay ventas de tiendas que no existen en el catálogo?
ventas |> anti_join(tiendas, by = "tienda_id") |> nrow()   # 0
\end{lstlisting}

\subsection{Fechas}

\begin{lstlisting}[language=R]
mensual <- ventas_det |>
  mutate(mes = floor_date(fecha, "month")) |>
  summarise(ingreso = sum(ingreso, na.rm = TRUE), .by = c(mes, canal)) |>
  arrange(canal, mes) |>
  group_by(canal) |>
  mutate(var_pct = (ingreso / lag(ingreso) - 1) * 100) |>
  ungroup()
mensual
\end{lstlisting}

\subsection{Texto}

\begin{lstlisting}[language=R]
ventas_det |>
  filter(str_detect(producto, "Smartphone")) |>
  mutate(producto = str_to_upper(producto)) |>
  distinct(producto)
\end{lstlisting}

\section{Actividades y ejercicios}

\begin{enumerate}
\def\labelenumi{\arabic{enumi}.}
\tightlist
\item
  ¿Cuál fue la \textbf{tienda} con mayor ingreso anual? Muestra el top 3
  con su canal y región.
\item
  Calcula la \textbf{participación (\%)} de cada canal en el ingreso
  total.
\item
  Para cada producto, encuentra la \textbf{semana} de mayores unidades
  vendidas (\passthrough{\lstinline!slice\_max()!}).
\item
  Compara el ingreso promedio semanal de \textbf{noviembre--diciembre}
  vs.~el resto del año por producto.
\item
  Reto: con \passthrough{\lstinline!anti\_join()!}, verifica que cada
  combinación tienda--semana tenga los 5 productos.
\end{enumerate}

\subsection{Soluciones}

\begin{lstlisting}[language=R]
# 1
ventas_det |>
  summarise(ingreso = sum(ingreso, na.rm = TRUE), .by = c(tienda_id, canal, region)) |>
  slice_max(ingreso, n = 3)

# 2
ventas_det |>
  summarise(ingreso = sum(ingreso, na.rm = TRUE), .by = canal) |>
  mutate(participacion = ingreso / sum(ingreso) * 100)

# 3
ventas |>
  summarise(unidades = sum(unidades, na.rm = TRUE), .by = c(producto, fecha)) |>
  group_by(producto) |>
  slice_max(unidades, n = 1, with_ties = FALSE)

# 4
ventas |>
  mutate(temporada = if_else(month(fecha) %in% 11:12, "Nov-Dic", "Resto")) |>
  summarise(ingreso_sem = sum(ingreso, na.rm = TRUE), .by = c(producto, temporada, fecha)) |>
  summarise(promedio = mean(ingreso_sem), .by = c(producto, temporada)) |>
  pivot_wider(names_from = temporada, values_from = promedio) |>
  mutate(incremento_pct = (`Nov-Dic` / Resto - 1) * 100)

# 5
esperado <- expand_grid(
  tienda_id = unique(ventas$tienda_id),
  fecha     = unique(ventas$fecha),
  producto  = unique(ventas$producto)
)
esperado |> anti_join(ventas, by = c("tienda_id", "fecha", "producto")) |> nrow()  # 0
\end{lstlisting}

\section{Evaluación (sumativa parcial)}

\textbf{Mini-reporte de KPIs:} script que produzca una tabla con
ingreso, unidades y ticket promedio por canal y región, más la variación
mensual por canal. Se evalúa: resultados correctos (40 \%), uso
idiomático de \passthrough{\lstinline!dplyr!} (30 \%), legibilidad (30
\%).

\section{Referencias verificadas}

\begin{itemize}
\tightlist
\item
  \textbf{Libro:} Wickham, H., Çetinkaya-Rundel, M., \& Grolemund, G.
  (2023). \emph{R for Data Science} (2.ª ed.). O'Reilly Media. ISBN
  978-1-4920-9740-2. Español: \url{https://davidrsch.github.io/r4ds-es/}
  {[}ES{]} {[}OA{]} --- caps. ``Transformación de datos'', ``Uniones'',
  ``Fechas y horas'', ``Cadenas''.
\item
  \textbf{Web:} R para Ciencia de Datos (traducción de la 1.ª ed.~por la
  comunidad latinoamericana de R). \url{https://es.r4ds.hadley.nz/}
  {[}ES{]} {[}OA{]} --- alternativa en español.
\item
  \textbf{Artículo:} Wickham, H., et al.~(2019). Welcome to the
  tidyverse. \emph{Journal of Open Source Software, 4}(43), 1686.
  \url{https://doi.org/10.21105/joss.01686} {[}EN{]} {[}OA{]}
\item
  \textbf{Video:} Gomila Salas, J. G. (Frogames). \emph{Curso completo
  de R para Data Science con Tidyverse} {[}Video{]}. YouTube.
  \url{https://www.youtube.com/watch?v=9NJyhs5PlGQ} {[}ES{]}
  \emph{{[}canal y fecha POR VERIFICAR{]}}
\end{itemize}

\chapter{Módulo 6 --- Visualización con ggplot2}

\begin{itemize}
\tightlist
\item
  \textbf{Duración estimada:} 1.5 h teoría + 2.5 h práctica
\item
  \textbf{Objetivos de aprendizaje.} Al finalizar, la persona podrá:

  \begin{enumerate}
  \def\labelenumi{\arabic{enumi}.}
  \tightlist
  \item
    \textbf{Explicar} la gramática de gráficos: datos, estéticas
    (\passthrough{\lstinline!aes!}), geometrías, escalas, facetas y
    temas.
  \item
    \textbf{Construir} gráficos de barras, líneas, dispersión,
    histogramas y cajas.
  \item
    \textbf{Seleccionar} el tipo de gráfico adecuado según la pregunta
    (comparar, tendencia, distribución, relación).
  \item
    \textbf{Producir} gráficos listos para reporte (títulos, etiquetas,
    formato de ejes) y exportarlos con
    \passthrough{\lstinline!ggsave()!}.
  \end{enumerate}
\end{itemize}

\section{Contenidos}

\begin{enumerate}
\def\labelenumi{\arabic{enumi}.}
\tightlist
\item
  La gramática de gráficos por capas.
\item
  \passthrough{\lstinline!aes()!}: \passthrough{\lstinline!x!},
  \passthrough{\lstinline!y!}, \passthrough{\lstinline!color!},
  \passthrough{\lstinline!fill!}, \passthrough{\lstinline!size!};
  estéticas fijas vs.~mapeadas.
\item
  Geometrías: \passthrough{\lstinline!geom\_col!},
  \passthrough{\lstinline!geom\_line!},
  \passthrough{\lstinline!geom\_point!},
  \passthrough{\lstinline!geom\_histogram!},
  \passthrough{\lstinline!geom\_boxplot!}.
\item
  Facetas: \passthrough{\lstinline!facet\_wrap()!}.
\item
  Escalas y formato: \passthrough{\lstinline!scales::label\_dollar()!},
  \passthrough{\lstinline!label\_percent()!}.
\item
  Etiquetas y temas: \passthrough{\lstinline!labs()!},
  \passthrough{\lstinline!theme\_minimal()!}.
\item
  Exportar: \passthrough{\lstinline!ggsave()!}.
\end{enumerate}

\section{Desarrollo}

\begin{lstlisting}[language=R]
library(tidyverse)
ventas  <- read_csv("datos/ventas.csv", show_col_types = FALSE) |>
  mutate(ingreso = precio * unidades)
tiendas <- read_csv("datos/tiendas.csv", show_col_types = FALSE)
ventas_det <- left_join(ventas, tiendas, by = "tienda_id")
\end{lstlisting}

\subsection{Comparar categorías: barras}

\begin{lstlisting}[language=R]
por_producto <- ventas |>
  summarise(ingreso = sum(ingreso, na.rm = TRUE), .by = producto)

ggplot(por_producto, aes(x = ingreso, y = fct_reorder(producto, ingreso))) +
  geom_col(fill = "steelblue") +
  scale_x_continuous(labels = scales::label_dollar(scale = 1e-6, suffix = " M")) +
  labs(title = "Ingreso anual por producto", x = "Ingreso (MXN)", y = NULL) +
  theme_minimal()
\end{lstlisting}

\begin{quote}
Barras \textbf{horizontales} y \textbf{ordenadas} facilitan leer nombres
largos y comparar magnitudes.
\end{quote}

\subsection{Tendencia: líneas}

\begin{lstlisting}[language=R]
semanal <- ventas_det |>
  summarise(ingreso = sum(ingreso, na.rm = TRUE), .by = c(fecha, canal))

g_tendencia <- ggplot(semanal, aes(fecha, ingreso, color = canal)) +
  geom_line(linewidth = 0.8) +
  scale_y_continuous(labels = scales::label_dollar()) +
  labs(title = "Ingreso semanal por canal, 2025", x = NULL, y = "Ingreso", color = "Canal") +
  theme_minimal()
g_tendencia
\end{lstlisting}

\subsection{Distribución: histograma y cajas}

\begin{lstlisting}[language=R]
ggplot(ventas, aes(unidades)) +
  geom_histogram(binwidth = 2, fill = "grey40", color = "white") +
  facet_wrap(~ producto, scales = "free_y") +
  labs(title = "Distribución de unidades semanales por producto")

ggplot(ventas_det, aes(region, unidades, fill = region)) +
  geom_boxplot(show.legend = FALSE) +
  labs(title = "Unidades por región", x = NULL)
\end{lstlisting}

\begin{quote}
Los \passthrough{\lstinline!NA!} de \passthrough{\lstinline!unidades!}
se descartan con una advertencia (``Removed rows\ldots{}''). Es
informativa, no un error.
\end{quote}

\subsection{Relación: dispersión}

\begin{lstlisting}[language=R]
por_tienda <- ventas_det |>
  summarise(unidades = sum(unidades, na.rm = TRUE),
            ingreso  = sum(ingreso,  na.rm = TRUE), .by = c(tienda_id, canal))

ggplot(por_tienda, aes(unidades, ingreso, color = canal)) +
  geom_point(size = 3) +
  geom_smooth(method = "lm", se = FALSE, color = "black", linewidth = 0.5) +
  labs(title = "Unidades vs. ingreso por tienda")
\end{lstlisting}

\subsection{Exportar}

\begin{lstlisting}[language=R]
ggsave("tendencia_canal.png", g_tendencia, width = 8, height = 4.5, dpi = 300)
\end{lstlisting}

\section{Actividades y ejercicios}

\begin{enumerate}
\def\labelenumi{\arabic{enumi}.}
\tightlist
\item
  Gráfico de barras apiladas al 100 \% de la participación de cada
  producto dentro de cada canal
  (\passthrough{\lstinline!position = "fill"!}).
\item
  Líneas de unidades mensuales por producto con
  \passthrough{\lstinline!facet\_wrap()!}.
\item
  Mejora un gráfico ``por defecto'' aplicando: orden, título que afirme
  la conclusión, ejes formateados y sin leyenda redundante.
\item
  Exporta tu mejor gráfico a PNG de 300 dpi.
\end{enumerate}

\subsection{Soluciones}

\begin{lstlisting}[language=R]
# 1
ventas_det |>
  summarise(ingreso = sum(ingreso, na.rm = TRUE), .by = c(canal, producto)) |>
  ggplot(aes(canal, ingreso, fill = producto)) +
  geom_col(position = "fill") +
  scale_y_continuous(labels = scales::label_percent()) +
  labs(title = "Mezcla de productos por canal", x = NULL, y = "Participación", fill = NULL)

# 2
ventas |>
  mutate(mes = floor_date(fecha, "month")) |>
  summarise(unidades = sum(unidades, na.rm = TRUE), .by = c(mes, producto)) |>
  ggplot(aes(mes, unidades)) +
  geom_line() +
  facet_wrap(~ producto, scales = "free_y") +
  labs(title = "Las ventas de todos los productos suben en noviembre y diciembre",
       x = NULL, y = "Unidades")
\end{lstlisting}

\section{Evaluación (formativa)}

Galería de pares: cada persona presenta un gráfico y recibe
retroalimentación con la lista de verificación: ¿responde una pregunta
clara?, ¿el tipo de gráfico es adecuado?, ¿ejes y unidades legibles?,
¿el título dice la conclusión?, ¿colores con propósito?

\section{Referencias verificadas}

\begin{itemize}
\tightlist
\item
  \textbf{Libro:} Wickham, H., Çetinkaya-Rundel, M., \& Grolemund, G.
  (2023). \emph{R for Data Science} (2.ª ed.). O'Reilly Media. ISBN
  978-1-4920-9740-2. Español: \url{https://davidrsch.github.io/r4ds-es/}
  {[}ES{]} {[}OA{]} --- caps. ``Visualización de datos'', ``Capas'',
  ``Análisis exploratorio'', ``Comunicación''.
\item
  \textbf{Web:} Wickham, H., Navarro, D., \& Pedersen, T. L.
  \emph{ggplot2: Elegant Graphics for Data Analysis} (3.ª ed., versión
  en línea). \url{https://ggplot2-book.org/} {[}EN{]} {[}OA en línea{]}
  \emph{{[}datos de edición impresa POR VERIFICAR{]}}
\item
  \textbf{Video:} freeCodeCamp.org. \emph{R Programming Tutorial --
  Learn the Basics of Statistical Computing} {[}Video{]}. YouTube.
  \url{https://youtu.be/_V8eKsto3Ug} {[}EN{]} --- secciones de gráficos
  (\passthrough{\lstinline!plot()!}, barras, histogramas, dispersión) en
  R base, útil para contrastar con ggplot2.
\end{itemize}

\chapter{Módulo 7 --- Estadística y modelos en R}

\begin{itemize}
\tightlist
\item
  \textbf{Duración estimada:} 1.5 h teoría + 2.5 h práctica
\item
  \textbf{Objetivos de aprendizaje.} Al finalizar, la persona podrá:

  \begin{enumerate}
  \def\labelenumi{\arabic{enumi}.}
  \tightlist
  \item
    \textbf{Calcular} estadísticos descriptivos (media, mediana,
    desviación estándar, cuantiles, correlación).
  \item
    \textbf{Ejecutar} pruebas de hipótesis básicas
    (\passthrough{\lstinline!t.test()!},
    \passthrough{\lstinline!wilcox.test()!},
    \passthrough{\lstinline!chisq.test()!}).
  \item
    \textbf{Ajustar} un modelo de regresión lineal con
    \passthrough{\lstinline!lm()!} usando la sintaxis de fórmulas.
  \item
    \textbf{Interpretar} coeficientes, valores p, R² y residuos, y
    \textbf{justificar} si el modelo es adecuado.
  \end{enumerate}
\end{itemize}

\section{Contenidos}

\begin{enumerate}
\def\labelenumi{\arabic{enumi}.}
\tightlist
\item
  Estadística descriptiva por grupo con \passthrough{\lstinline!dplyr!}.
\item
  Correlación: \passthrough{\lstinline!cor()!} (Pearson y Spearman).
\item
  Pruebas de hipótesis: t de Welch, Mann-Whitney/Wilcoxon, chi-cuadrada.
\item
  Fórmulas en R: \passthrough{\lstinline!y \~ x1 + x2!}.
\item
  Regresión lineal: \passthrough{\lstinline!lm()!},
  \passthrough{\lstinline!summary()!}, \passthrough{\lstinline!coef()!},
  \passthrough{\lstinline!confint()!},
  \passthrough{\lstinline!predict()!}.
\item
  Diagnóstico con residuos; resultados ordenados con
  \passthrough{\lstinline!broom::tidy()!}.
\end{enumerate}

\section{Desarrollo}

\begin{lstlisting}[language=R]
library(tidyverse)
ventas  <- read_csv("datos/ventas.csv", show_col_types = FALSE)
tiendas <- read_csv("datos/tiendas.csv", show_col_types = FALSE)
ventas_det <- ventas |>
  left_join(tiendas, by = "tienda_id") |>
  drop_na(unidades) |>
  mutate(temporada = if_else(month(fecha) %in% 11:12, "Nov-Dic", "Resto"))
\end{lstlisting}

\subsection{Descriptivos}

\begin{lstlisting}[language=R]
ventas_det |>
  summarise(
    media   = mean(unidades),
    mediana = median(unidades),
    de      = sd(unidades),
    p90     = quantile(unidades, 0.9),
    .by = producto
  )
\end{lstlisting}

\subsection{Prueba t: ¿la temporada cambia las ventas?}

\begin{lstlisting}[language=R]
reloj <- ventas_det |> filter(producto == "Reloj")
t.test(unidades ~ temporada, data = reloj)       # Welch por defecto
wilcox.test(unidades ~ temporada, data = reloj)  # alternativa no paramétrica
\end{lstlisting}

\begin{quote}
\textbf{Cuidado:} las observaciones semanales de una misma tienda no son
del todo independientes. Aquí la prueba es ilustrativa; en un análisis
formal hay que considerar la autocorrelación.
\end{quote}

\subsection{Chi-cuadrada: ¿la mezcla de productos depende del canal?}

\begin{lstlisting}[language=R]
tabla <- xtabs(unidades ~ canal + producto, data = ventas_det)
tabla
chisq.test(tabla)
\end{lstlisting}

\subsection{Regresión lineal}

\begin{lstlisting}[language=R]
semanal <- ventas_det |>
  summarise(unidades = sum(unidades), .by = c(fecha, canal, temporada))

modelo <- lm(unidades ~ temporada + canal, data = semanal)
summary(modelo)
confint(modelo)

library(broom)
tidy(modelo)          # coeficientes como tibble
glance(modelo)        # R², AIC, etc.

nuevo <- tibble(temporada = "Nov-Dic", canal = "En línea")
predict(modelo, newdata = nuevo, interval = "prediction")
\end{lstlisting}

\subsection{Diagnóstico}

\begin{lstlisting}[language=R]
aug <- augment(modelo)
ggplot(aug, aes(.fitted, .resid)) +
  geom_point(alpha = 0.5) +
  geom_hline(yintercept = 0, linetype = "dashed") +
  labs(title = "Residuos vs. ajustados", x = "Ajustado", y = "Residuo")
\end{lstlisting}

\section{Actividades y ejercicios}

\begin{enumerate}
\def\labelenumi{\arabic{enumi}.}
\tightlist
\item
  Calcula la matriz de correlación entre
  \passthrough{\lstinline!precio!} y \passthrough{\lstinline!unidades!}
  a nivel registro. ¿Qué signo esperas y por qué?
\item
  Prueba si las unidades de \textbf{Audífonos} difieren entre las
  regiones \textbf{Norte} y \textbf{Sur}.
\item
  Ajusta \passthrough{\lstinline!lm(unidades \~ temporada * canal)!} y
  compara su R² ajustado con el modelo sin interacción.
\item
  Redacta en 3 líneas la interpretación del coeficiente
  \passthrough{\lstinline!temporadaResto!} para un público no técnico.
\end{enumerate}

\subsection{Soluciones}

\begin{lstlisting}[language=R]
# 1: negativo; los productos caros se venden en menos unidades
cor(ventas_det$precio, ventas_det$unidades)
cor(ventas_det$precio, ventas_det$unidades, method = "spearman")

# 2
audif <- ventas_det |> filter(producto == "Audífonos", region %in% c("Norte", "Sur"))
t.test(unidades ~ region, data = audif)

# 3
modelo_int <- lm(unidades ~ temporada * canal, data = semanal)
c(sin_interaccion = glance(modelo)$adj.r.squared,
  con_interaccion = glance(modelo_int)$adj.r.squared)
anova(modelo, modelo_int)

# 4 (ejemplo): "Fuera de noviembre-diciembre, cada canal vende en promedio
# alrededor de X unidades menos por semana que en temporada alta, manteniendo
# el canal constante; la diferencia es estadísticamente clara (p < 0.001)."
coef(modelo)["temporadaResto"]
\end{lstlisting}

\section{Evaluación (sumativa parcial)}

Informe breve (1 página) con: pregunta de negocio, prueba o modelo
elegido y \textbf{por qué}, resultado con intervalo de confianza,
verificación de supuestos y una limitación. Rúbrica: pertinencia del
método (30 \%), ejecución correcta (30 \%), interpretación (30 \%),
comunicación (10 \%).

\section{Referencias verificadas}

\begin{itemize}
\tightlist
\item
  \textbf{Video:} StatQuest with Josh Starmer. \emph{Linear Regression
  in R, Step-by-Step} {[}Video{]}. YouTube.
  \url{https://www.youtube.com/watch?v=u1cc1r_Y7M0} {[}EN{]}
  \emph{{[}fecha POR VERIFICAR{]}}
\item
  \textbf{Video:} StatQuest with Josh Starmer. \emph{Linear Regression
  and Linear Models} {[}Lista de reproducción{]}. YouTube.
  \url{https://www.youtube.com/playlist?list=PLblh5JKOoLUIzaEkCLIUxQFjPIlapw8nU}
  {[}EN{]}
\item
  \textbf{Web:} R Core Team. \emph{An Introduction to R} --- sección
  ``Statistical models in R''.
  \url{https://cran.r-project.org/doc/manuals/r-release/R-intro.html}
  {[}EN{]} {[}OA{]}
\item
  \textbf{Libro:} Wickham, H., Çetinkaya-Rundel, M., \& Grolemund, G.
  (2023). \emph{R for Data Science} (2.ª ed.). O'Reilly Media. ISBN
  978-1-4920-9740-2. Español: \url{https://davidrsch.github.io/r4ds-es/}
  {[}ES{]} {[}OA{]} --- caps. de análisis exploratorio como base para
  modelar.
\end{itemize}

\chapter{Módulo 8 --- Reportes reproducibles con Quarto y proyecto
integrador}

\begin{itemize}
\tightlist
\item
  \textbf{Duración estimada:} 1 h teoría + 4 h proyecto
\item
  \textbf{Objetivos de aprendizaje.} Al finalizar, la persona podrá:

  \begin{enumerate}
  \def\labelenumi{\arabic{enumi}.}
  \tightlist
  \item
    \textbf{Crear} un documento Quarto (\passthrough{\lstinline!.qmd!})
    que combine texto, código R y resultados.
  \item
    \textbf{Configurar} opciones de \emph{chunks}
    (\passthrough{\lstinline!echo!}, \passthrough{\lstinline!warning!},
    \passthrough{\lstinline!fig-cap!}) y formatos de salida (HTML, Word,
    PDF).
  \item
    \textbf{Organizar} un proyecto reproducible (estructura de carpetas,
    rutas relativas, \passthrough{\lstinline!renv!}).
  \item
    \textbf{Diseñar} y \textbf{producir} un análisis completo de
    principio a fin (proyecto integrador).
  \end{enumerate}
\end{itemize}

\section{Contenidos}

\begin{enumerate}
\def\labelenumi{\arabic{enumi}.}
\tightlist
\item
  Por qué la reproducibilidad: el código \emph{es} la documentación del
  análisis.
\item
  Anatomía de un \passthrough{\lstinline!.qmd!}: encabezado YAML,
  Markdown y \emph{chunks} de código.
\item
  Opciones de \emph{chunk} con \passthrough{\lstinline!\#|!}.
\item
  Tablas con \passthrough{\lstinline!knitr::kable()!}; gráficos con
  leyenda.
\item
  Renderizar: botón \textbf{Render} en RStudio o
  \passthrough{\lstinline!quarto render!} en terminal.
\item
  Estructura de proyecto y control de dependencias con
  \passthrough{\lstinline!renv!}.
\item
  Guía de estilo del tidyverse; control de versiones con Git
  (introducción).
\end{enumerate}

\section{Desarrollo}

\subsection{Un documento Quarto mínimo}

Guarda esto como \passthrough{\lstinline!reporte.qmd!} dentro del
proyecto y presiona \textbf{Render}:

\begin{lstlisting}
---
title: "Reporte de ventas 2025"
author: "Tu nombre"
format:
  html:
    toc: true
execute:
  warning: false
---

## Resumen

```{r}
#| label: datos
#| echo: false
library(tidyverse)
ventas <- read_csv("datos/ventas.csv", show_col_types = FALSE) |>
  mutate(ingreso = precio * unidades)
total <- sum(ventas$ingreso, na.rm = TRUE)
```

El ingreso total del año fue de **`r scales::dollar(total)`**.

```{r}
#| label: fig-producto
#| fig-cap: "Ingreso anual por producto"
ventas |>
  summarise(ingreso = sum(ingreso, na.rm = TRUE), .by = producto) |>
  ggplot(aes(ingreso, fct_reorder(producto, ingreso))) +
  geom_col() +
  labs(x = "Ingreso", y = NULL)
```
\end{lstlisting}

\begin{itemize}
\tightlist
\item
  \passthrough{\lstinline!`r expresion`!} inserta resultados
  \textbf{dentro del texto}: si cambian los datos, el número se
  actualiza solo.
\item
  Cambia \passthrough{\lstinline!format: html!} por
  \passthrough{\lstinline!docx!} o \passthrough{\lstinline!pdf!} (este
  último requiere TinyTeX:
  \passthrough{\lstinline!quarto install tinytex!}).
\end{itemize}

\subsection{Proyecto reproducible}

\begin{lstlisting}
mi-analisis/
|-- mi-analisis.Rproj
|-- datos/          # datos crudos: nunca se editan a mano
|-- R/              # funciones reutilizables
|-- reporte.qmd
`-- salidas/        # gráficos y tablas generadas
\end{lstlisting}

\begin{lstlisting}[language=R]
# Registrar versiones de paquetes del proyecto (una vez, en la consola):
# install.packages("renv"); renv::init()
# Después de instalar o actualizar paquetes:
# renv::snapshot()
\end{lstlisting}

\section{Proyecto integrador (capstone)}

\textbf{Caso:} Eres analista de una cadena de electrónica. Dirección
pide un \textbf{reporte de desempeño 2025} con recomendaciones para la
temporada alta de 2026.

\textbf{Entregables:}

\begin{enumerate}
\def\labelenumi{\arabic{enumi}.}
\tightlist
\item
  Proyecto de RStudio con la estructura anterior.
\item
  \passthrough{\lstinline!reporte.qmd!} renderizado a HTML que incluya:

  \begin{itemize}
  \tightlist
  \item
    Limpieza documentada (faltantes, tipos, verificación de llaves).
  \item
    Tabla de KPIs por canal y región (ingreso, unidades, ticket
    promedio, participación).
  \item
    Al menos 3 gráficos con títulos que afirmen la conclusión.
  \item
    Una prueba estadística o modelo con interpretación.
  \item
    Al menos una función propia en \passthrough{\lstinline!R/!} usada en
    el reporte.
  \item
    Conclusiones y 3 recomendaciones accionables.
  \end{itemize}
\item
  Presentación oral de 5 minutos.
\end{enumerate}

\subsection{Rúbrica global del proyecto}

\begin{longtable}[]{@{}
  >{\raggedright\arraybackslash}p{(\columnwidth - 8\tabcolsep) * \real{0.2000}}
  >{\raggedright\arraybackslash}p{(\columnwidth - 8\tabcolsep) * \real{0.2000}}
  >{\raggedright\arraybackslash}p{(\columnwidth - 8\tabcolsep) * \real{0.2000}}
  >{\raggedright\arraybackslash}p{(\columnwidth - 8\tabcolsep) * \real{0.2000}}
  >{\raggedright\arraybackslash}p{(\columnwidth - 8\tabcolsep) * \real{0.2000}}@{}}
\toprule\noalign{}
\begin{minipage}[b]{\linewidth}\raggedright
Criterio (peso)
\end{minipage} & \begin{minipage}[b]{\linewidth}\raggedright
4 --- Excelente
\end{minipage} & \begin{minipage}[b]{\linewidth}\raggedright
3 --- Bueno
\end{minipage} & \begin{minipage}[b]{\linewidth}\raggedright
2 --- Suficiente
\end{minipage} & \begin{minipage}[b]{\linewidth}\raggedright
1 --- Insuficiente
\end{minipage} \\
\midrule\noalign{}
\endhead
\bottomrule\noalign{}
\endlastfoot
Reproducibilidad (20 \%) & Renderiza sin errores en otra computadora;
rutas relativas; \passthrough{\lstinline!renv!} & Renderiza con ajustes
menores & Requiere varios ajustes & No renderiza \\
Manejo de datos (20 \%) & Limpieza completa y justificada; verificación
de llaves & Limpieza correcta, poco documentada & Limpieza parcial &
Datos sin revisar \\
Transformación y KPIs (20 \%) & KPIs correctos, código
\passthrough{\lstinline!dplyr!} idiomático & Correctos con código
redundante & Errores menores & Errores graves \\
Visualización (15 \%) & Gráficos adecuados, claros y con mensaje &
Adecuados, mejorables & Tipo de gráfico poco apropiado & Ilegibles \\
Estadística (15 \%) & Método pertinente, supuestos revisados,
interpretación correcta & Pertinente, interpretación parcial & Método
dudoso & Ausente o incorrecto \\
Comunicación (10 \%) & Conclusiones y recomendaciones claras y
accionables & Claras pero generales & Confusas & Ausentes \\
\end{longtable}

\section{Referencias verificadas}

\begin{itemize}
\tightlist
\item
  \textbf{Web:} Posit / Quarto. \emph{Tutorial: Computations} (RStudio).
  \url{https://quarto.org/docs/get-started/computations/rstudio}
  {[}EN{]} {[}gratuito{]}
\item
  \textbf{Web:} Posit / Quarto. \emph{Using R}.
  \url{https://quarto.org/docs/computations/r.html} {[}EN{]}
  {[}gratuito{]}
\item
  \textbf{Web:} Posit / Quarto. \emph{Tutorial: Hello, Quarto}.
  \url{https://quarto.org/docs/get-started/hello/rstudio.html} {[}EN{]}
  {[}gratuito{]}
\item
  \textbf{Libro:} Wickham, H., Çetinkaya-Rundel, M., \& Grolemund, G.
  (2023). \emph{R for Data Science} (2.ª ed.). O'Reilly Media. ISBN
  978-1-4920-9740-2. Español: \url{https://davidrsch.github.io/r4ds-es/}
  {[}ES{]} {[}OA{]} --- caps. ``Quarto'' y ``Formatos de Quarto''.
\item
  \textbf{Artículo:} Wickham, H. (2014). Tidy data. \emph{Journal of
  Statistical Software, 59}(10), 1--23.
  \url{https://doi.org/10.18637/jss.v059.i10} {[}EN{]} {[}OA{]}
\end{itemize}


\end{document}
